% ====================================================================
% ФАЙЛ: tasks/03_electrodynamics/electric_potential_part1.tex
% ТЕМА: ПОТЕНЦИАЛ ЭЛЕКТРИЧЕСКОГО ПОЛЯ
% ПАЧКА 1: ВАРИАНТЫ 1–10
% ====================================================================

% === ЗАДАЧА 1 ===
% Тег: #ЭЛЕКТРОДИНАМИКА #ПОТЕНЦИАЛ #КИМ_11 #ВАРИАНТ_01
\begin{taskbasic}[code={\label{task:elem_pot_v1_n11}}]
\textbf{Задача 1.} Какую работу совершают силы однородного электрического поля напряжённостью $E = 500\text{ В/м}$ при перемещении точечного заряда $q = 2\text{ мкКл}$ на расстояние $d = 4\text{ см}$ вдоль силовой линии поля? Ответ выразите в микроджоулях (мкДж).

\kim{11}
\end{taskbasic}
\answer{40}

% === ЗАДАЧА 2 ===
% Тег: #ЭЛЕКТРОДИНАМИКА #ПОТЕНЦИАЛ #КИМ_11 #ВАРИАНТ_02
\begin{taskbasic}[code={\label{task:elem_pot_v2_n11}}]
\textbf{Задача 2.} Точечный заряд $q = 1\text{ мкКл}$ перемещается в электростатическом поле из точки с потенциалом $\varphi_1 = 50\text{ В}$ в точку с потенциалом $\varphi_2 = 30\text{ В}$. Какую работу совершают силы электрического поля при этом перемещении? Ответ выразите в микроджоулях (мкДж).

\kim{11}
\end{taskbasic}
\answer{20}

% === ЗАДАЧА 3 ===
% Тег: #ЭЛЕКТРОДИНАМИКА #ПОТЕНЦИАЛ #КИМ_15 #ВАРИАНТ_05
\begin{taskbasic}[code={\label{task:elem_pot_v5_n15}}]
\textbf{Задача 3.} На рисунке показан график зависимости потенциала $\varphi$ электростатического поля, создаваемого некоторой системой зарядов, от координаты $x$. Из приведённого ниже списка выберите все верные утверждения относительно характеристик этого поля.

\nopagebreak\smallskip
\begin{center}
\begin{tikzpicture}[scale=1.0, every node/.style={font=\footnotesize}]
    \draw[xstep=1, ystep=1, gray!30, thin] (0,0) grid (6,4);
    \draw[-{Stealth[scale=0.8]}, black, thick] (0,0) -- (6.5,0) node[right] {$x$, см};
    \draw[-{Stealth[scale=0.8]}, black, thick] (0,0) -- (0,4.5) node[above] {$\varphi$, В};
    
    \coordinate (A) at (0,3);
    \coordinate (B) at (2,3);
    \coordinate (C) at (4,1);
    \coordinate (D) at (6,1);
    
    \draw[ultra thick, brandPrimary] (A) -- (B) -- (C) -- (D);
    
    \draw[dashed, gray] (0,3) -- (B);
    \draw[dashed, gray] (0,1) -- (C);
    \draw[dashed, gray] (2,0) -- (B);
    \draw[dashed, gray] (4,0) -- (C);
    \draw[dashed, gray] (6,0) -- (D);
    
    \node[left] at (0,1) {10};
    \node[left] at (0,3) {30};
    \node[below] at (2,0) {2};
    \node[below] at (4,0) {4};
    \node[below] at (6,0) {6};
    \node[below left] at (0,0) {0};
\end{tikzpicture}
\end{center}

1) На участке от $x = 0$ до $x = 2\text{ см}$ проекция напряжённости поля $E_x = 0$.\\
2) Проекция напряжённости поля $E_x$ на участке от $x = 2\text{ см}$ до $x = 4\text{ см}$ равна $-10\text{ В/см}$.\\
3) В точке $x = 3\text{ см}$ вектор напряжённости поля направлен вдоль положительного направления оси $Ox$.\\
4) Модуль напряжённости поля на участке от $x = 4\text{ см}$ до $x = 6\text{ см}$ максимален.\\
5) На участке от $x = 2\text{ см}$ до $x = 4\text{ см}$ потенциальная энергия положительного заряда уменьшается.

\kim{15}
\end{taskbasic}
\answer{13}

% === ЗАДАЧА 4 ===
% Тег: #ЭЛЕКТРОДИНАМИКА #ПОТЕНЦИАЛ #КИМ_15 #ВАРИАНТ_06
\begin{taskbasic}[code={\label{task:elem_pot_v6_n15}}]
\textbf{Задача 4.} На рисунке показан график зависимости потенциала $\varphi$ электростатического поля от координаты $x$. Из приведённого ниже списка выберите все верные утверждения относительно характеристик этого поля.

\nopagebreak\smallskip
\begin{center}
\begin{tikzpicture}[scale=1.0, every node/.style={font=\footnotesize}]
    \draw[xstep=1, ystep=1, gray!30, thin] (0,0) grid (6,4);
    \draw[-{Stealth[scale=0.8]}, black, thick] (0,0) -- (6.5,0) node[right] {$x$, см};
    \draw[-{Stealth[scale=0.8]}, black, thick] (0,0) -- (0,4.5) node[above] {$\varphi$, В};
    
    \coordinate (A) at (0,1);
    \coordinate (B) at (2,3);
    \coordinate (C) at (4,3);
    \coordinate (D) at (6,0);
    
    \draw[ultra thick, brandPrimary] (A) -- (B) -- (C) -- (D);
    
    \draw[dashed, gray] (0,1) -- (A);
    \draw[dashed, gray] (0,3) -- (B);
    \draw[dashed, gray] (2,0) -- (B);
    \draw[dashed, gray] (4,0) -- (C);
    \draw[dashed, gray] (6,0) -- (D);
    
    \node[left] at (0,1) {10};
    \node[left] at (0,3) {30};
    \node[below] at (2,0) {2};
    \node[below] at (4,0) {4};
    \node[below] at (6,0) {6};
    \node[below left] at (0,0) {0};
\end{tikzpicture}
\end{center}

1) На участке от $x = 0$ до $x = 2\text{ см}$ модуль напряжённости поля равен $20\text{ В/м}$.\\
2) На участке от $x = 2\text{ см}$ до $x = 4\text{ см}$ напряжённость электрического поля равна нулю.\\
3) В точке $x = 1\text{ см}$ вектор напряжённости поля направлен вдоль положительного направления оси $Ox$.\\
4) В точке $x = 5\text{ см}$ вектор напряжённости поля направлен в сторону положительного направления оси $Ox$.\\
5) Работа поля при перемещении положительного заряда из точки $x = 2\text{ см}$ в точку $x = 4\text{ см}$ отрицательна.

\kim{15}
\end{taskbasic}
\answer{24}

% ====================================================================
% ФАЙЛ: tasks/03_electrodynamics/electric_potential_part2.tex
% ТЕМА: ПОТЕНЦИАЛ ЭЛЕКТРИЧЕСКОГО ПОЛЯ
% ПАЧКА 2: ВАРИАНТЫ 11–20
% ====================================================================

% === ЗАДАЧА 1 ===
% Тег: #ЭЛЕКТРОДИНАМИКА #ПОТЕНЦИАЛ #КИМ_11 #ВАРИАНТ_13
\begin{taskbasic}[code={\label{task:elem_pot_v13_n11}}]
\textbf{Задача 1.} Точечный заряд $q = 3\text{ мкКл}$ перемещается в электростатическом поле из точки с потенциалом $\varphi_1 = 30\text{ В}$ в точку с потенциалом $\varphi_2 = 20\text{ В}$. Какую работу совершают силы электрического поля при этом перемещении? Ответ выразите в микроджоулях (мкДж).

\kim{11}
\end{taskbasic}
\answer{30}

% === ЗАДАЧА 2 ===
% Тег: #ЭЛЕКТРОДИНАМИКА #ПОТЕНЦИАЛ #КИМ_11 #ВАРИАНТ_14
\begin{taskbasic}[code={\label{task:elem_pot_v14_n11}}]
\textbf{Задача 2.} Какую работу совершают силы однородного электростатического поля при перемещении заряда $q = 5\text{ мкКл}$ между точками, разность потенциалов между которыми равна $30\text{ В}$? Ответ выразите в микроджоулях (мкДж).

\kim{11}
\end{taskbasic}
\answer{150}

% === ЗАДАЧА 3 ===
% Тег: #ЭЛЕКТРОДИНАМИКА #ПОТЕНЦИАЛ #КИМ_15 #ВАРИАНТ_17
\begin{taskbasic}[code={\label{task:elem_pot_v17_n15}}]
\textbf{Задача 3.} На рисунке изображены линии напряжённости однородного электростатического поля и четыре эквипотенциальные поверхности, потенциалы которых равны $10\text{ В}$, $20\text{ В}$, $30\text{ В}$ и $40\text{ В}$. Расстояние между соседними эквипотенциальными поверхностями равно $5\text{ см}$. Из приведённого ниже списка выберите все верные утверждения.

\nopagebreak\smallskip
\begin{center}
\begin{tikzpicture}[scale=1.0, every node/.style={font=\footnotesize}]
    \draw[xstep=1, ystep=1, gray!30, thin] (0,0) grid (4,3);
    
    % Эквипотенциальные линии (вертикальные)
    \draw[dashed, gray!80, thick] (1,0) -- (1,3);
    \draw[dashed, gray!80, thick] (2,0) -- (2,3);
    \draw[dashed, gray!80, thick] (3,0) -- (3,3);
    \draw[dashed, gray!80, thick] (4,0) -- (4,3);
    
    % Силовые линии (горизонтальные, влево)
    \draw[-{Stealth[scale=0.8]}, brandPrimary, thick] (4,0.8) -- (0.5,0.8);
    \draw[-{Stealth[scale=0.8]}, brandPrimary, thick] (4,2.2) -- (0.5,2.2) node[above right] {$\vec{E}$};
    
    % Подписи потенциалов
    \node[above] at (1,3) {$10\text{ В}$};
    \node[above] at (2,3) {$20\text{ В}$};
    \node[above] at (3,3) {$30\text{ В}$};
    \node[above] at (4,3) {$40\text{ В}$};
    
    \node[below left] at (0,0) {0};
\end{tikzpicture}
\end{center}

1) Модуль напряжённости электрического поля равен $2\text{ В/м}$.\\
2) Модуль напряжённости электрического поля равен $200\text{ В/м}$.\\
3) Вектор напряжённости поля направлен слева направо.\\
4) Работа поля по перемещению положительного заряда из точки с потенциалом $40\text{ В}$ в точку с потенциалом $10\text{ В}$ отрицательна.\\
5) Работа сил поля по перемещению заряда вдоль одной эквипотенциальной поверхности равна нулю.

\kim{15}
\end{taskbasic}
\answer{25}

% ====================================================================
% ФАЙЛ: tasks/03_electrodynamics/electric_potential_part3.tex
% ТЕМА: ПОТЕНЦИАЛ ЭЛЕКТРИЧЕСКОГО ПОЛЯ
% ПАЧКА 3: ВАРИАНТЫ 21–30
% ====================================================================

% === ЗАДАЧА 1 ===
% Тег: #ЭЛЕКТРОДИНАМИКА #ПОТЕНЦИАЛ #КИМ_11 #ВАРИАНТ_23
\begin{taskbasic}[code={\label{task:elem_pot_v23_n11}}]
\textbf{Задача 1.} Какую работу совершают силы однородного электростатического поля при перемещении заряда $q = 4\text{ мкКл}$ между точками, разность потенциалов между которыми равна $30\text{ В}$? Ответ выразите в микроджоулях (мкДж).

\kim{11}
\end{taskbasic}
\answer{120}

% === ЗАДАЧА 2 ===
% Тег: #ЭЛЕКТРОДИНАМИКА #ПОТЕНЦИАЛ #КИМ_11 #ВАРИАНТ_24
\begin{taskbasic}[code={\label{task:elem_pot_v24_n11}}]
\textbf{Задача 2.} Точечный заряд $q = 2\text{ мкКл}$ перемещается в электростатическом поле из точки с потенциалом $\varphi_1 = 60\text{ В}$ в точку с потенциалом $\varphi_2 = 40\text{ В}$. Какую работу совершают силы электрического поля при этом перемещении? Ответ выразите в микроджоулях (мкДж).

\kim{11}
\end{taskbasic}
\answer{40}

% === ЗАДАЧА 3 ===
% Тег: #ЭЛЕКТРОДИНАМИКА #ПОТЕНЦИАЛ #КИМ_15 #ВАРИАНТ_27
\begin{taskbasic}[code={\label{task:elem_pot_v27_n15}}]
\textbf{Задача 3.} На рисунке показан график зависимости потенциала $\varphi$ электростатического поля от координаты $x$. Из приведённого ниже списка выберите все верные утверждения относительно характеристик этого поля.

\nopagebreak\smallskip
\begin{center}
\begin{tikzpicture}[scale=1.0, every node/.style={font=\footnotesize}]
    \draw[xstep=1, ystep=1, gray!30, thin] (0,0) grid (6,4);
    \draw[-{Stealth[scale=0.8]}, black, thick] (0,0) -- (6.5,0) node[right] {$x$, см};
    \draw[-{Stealth[scale=0.8]}, black, thick] (0,0) -- (0,4.5) node[above] {$\varphi$, В};
    
    \coordinate (A) at (0,4);
    \coordinate (B) at (2,2);
    \coordinate (C) at (4,2);
    \coordinate (D) at (6,0);
    
    \draw[ultra thick, brandPrimary] (A) -- (B) -- (C) -- (D);
    
    \draw[dashed, gray] (0,4) -- (A);
    \draw[dashed, gray] (0,2) -- (B);
    \draw[dashed, gray] (2,0) -- (B);
    \draw[dashed, gray] (4,0) -- (C);
    \draw[dashed, gray] (6,0) -- (D);
    
    \node[left] at (0,2) {20};
    \node[left] at (0,4) {40};
    \node[below] at (2,0) {2};
    \node[below] at (4,0) {4};
    \node[below] at (6,0) {6};
    \node[below left] at (0,0) {0};
\end{tikzpicture}
\end{center}

1) На участке от $x = 0$ до $x = 2\text{ см}$ модуль напряжённости поля равен нулю.\\
2) На участке от $x = 2\text{ см}$ до $x = 4\text{ см}$ напряжённость электрического поля равна нулю.\\
3) Проекция напряжённости поля $E_x$ на участке от $x = 0$ до $x = 2\text{ см}$ положительна и равна $10\text{ В/см}$.\\
4) Модуль напряжённости поля в точке $x = 5\text{ см}$ равен $1000\text{ В/м}$.\\
5) Работа поля при перемещении положительного заряда из точки $x = 2\text{ см}$ в точку $x = 4\text{ см}$ положительна.

\kim{15}
\end{taskbasic}
\answer{23}

% === ЗАДАЧА 4 ===
% Тег: #ЭЛЕКТРОДИНАМИКА #ПОТЕНЦИАЛ #КИМ_25 #ВАРИАНТ_28
\begin{taskbasic}[code={\label{task:elem_pot_v28_n25}}]
\textbf{Задача 4.} Пылинка массой $m = 10^{-8}\text{ г}$ и зарядом $q = 5\cdot 10^{-13}\text{ Кл}$ влетает в однородное электрическое поле вдоль силовой линии со скоростью $v_0 = 3\text{ м/с}$. Какова напряжённость электрического поля $E$, если до остановки пылинка проходит расстояние $d = 2\text{ см}$? Ответ выразите в В/м.

\kim{25}
\end{taskbasic}
\answer{600}